% !TEX program = xelatex
% =============================================================================
%  中国石油大学（北京）本科生毕业设计(论文)
%  由 md2cupb.py 自动生成
%  编译: xelatex → bibtex → xelatex → xelatex
% =============================================================================

\documentclass{cupb-thesis}

% ---------------------------------------------------------------------------
%  论文元信息 (请根据实际情况修改)
% ---------------------------------------------------------------------------
\thesititle{中国石油大学论文写作指南}
\thesiauthor{教务处}
\thesistudentid{000000}
\thesicollege{教务处}
\thesimajor{教学管理}
\thesiadviser{指导教师}
\thesifinishdate{2026年5月}
\thesientitle{English Title}
\cnkeywords{关键词1；关键词2；关键词3}
\enkeywords{keyword1; keyword2; keyword3}

% ---------------------------------------------------------------------------
%  参考文献文件 (如果有 .bib 文件, 在此指定)
% ---------------------------------------------------------------------------
% \addbibresource{references.bib}

\begin{document}

% ===========================================================================
%  前置部分
% ===========================================================================

%% --- 中文封面 ---
\makecover

%% --- 声明页 ---
\makestatement

%% --- 中文摘要 ---
\begin{cnabstract}
（请在此处填写中文摘要，300-500字）
\end{cnabstract}

%% --- 英文摘要 ---
\begin{enabstract}
(Please write the English abstract here, 300-500 words.)
\end{enabstract}

%% --- 目录 ---
\tableofcontents

%% --- 前言 ---
\begin{preface}
（请在此处填写前言内容）
\end{preface}

% ===========================================================================
%  正文 (从第一章开始, 页码用阿拉伯数字)
% ===========================================================================
\mainmatterstart
% 文档标题: 本科生毕业设计（论文）写作指南

本科生毕业设计(论文)是实现人才培养目标的重要实践性环节，对巩固、深化和升华学生所学理论知识，培养学生创新能力、独立工作能力、分析和解决问题能力、工程实践能力起着重要作用。毕业设计(论文)同时也是记录科研成果的重要文献资料，是申请学位的基本依据。为了规范我校本科生毕业设计(论文)的管理，保证毕业设计

(论文)质量，特制定《中国石油大学（北京）本科生毕业设计（论文）写作指南》。

\section{一、本科生毕业设计(论文)应体现的三个规范}

本科生毕业设计(论文)的形式与格式虽然不能直接反映论文的学术水平，但可体现出作者的学术修养和文化修养。合格的本科生毕业设计(论文)在形式与格式上体现如下三方面规范：

\begin{enumerate}[label=\arabic*.]
  \item 国家学位条例和学校对学位论文管理的各项具体规范。
  \item 学术著述的一般规范。
  \item 所有正式出版物共同遵循的文字印刷规范。
\end{enumerate}
\section{二、本科生毕业设计(论文)和摘要篇幅要求}

\begin{enumerate}[label=\arabic*.]
  \item 本科生毕业设计(论文)一般为1.5万字左右。
  \item 中文、英文摘要一般为300-500字。
\end{enumerate}
\section{三、本科生毕业设计(论文)封面书写要求}

论文封面包括以下9项：

\begin{enumerate}[label=（\arabic*）]
  \item 单位代码    中国石油大学（北京）的代码是11414；
  \item 学号        以《本科生名册》中的本科生学号为准；
  \item 论文标志    “中国石油大学本科生毕业
\end{enumerate}
\section{一、本科生毕业设计(论文)应体现的三个规范}

本科生毕业设计(论文)的形式与格式虽然不能直接反映论文的学术水平，但可体现出作者的学术修养和文化修养。合格的本科生毕业设计(论文)在形式与格式上体现如下三方面规范：

\begin{enumerate}[label=\arabic*.]
  \item 国家学位条例和学校对学位论文管理的各项具体规范。
  \item 学术著述的一般规范。
  \item 所有正式出版物共同遵循的文字印刷规范。
\end{enumerate}
\section{二、本科生毕业设计(论文)和摘要篇幅要求}

\begin{enumerate}[label=\arabic*.]
  \item 本科生毕业设计(论文)一般为1.5万字左右。
  \item 中文、英文摘要一般为300-500字。
\end{enumerate}
\section{三、本科生毕业设计(论文)封面书写要求}

论文封面包括以下9项：

\begin{enumerate}[label=（\arabic*）]
  \item 单位代码    中国石油大学（北京）的代码是11414；
  \item 学号        以《本科生名册》中的本科生学号为准；
  \item 论文标志    “中国石油大学本科生毕业设计(论文)”；
  \item 论文题目    应简短、醒目、切题，不宜超过25个字；
  \item 学院名称以《本科生培养计划》中的学院为准；
  \item 专业名称以《本科生培养计划》中的专业名称为准；
  \item 学生姓名    以《本科生名册》中的姓名为准；
  \item 指导教师    论文的指导教师；
  \item 论文完成时间    指论文起止时间，用阿拉伯数字书写。
\end{enumerate}
\section{四、本科生毕业设计(论文)摘要书写内容及\# 本科生毕业设计（论文）写作指南}

本科生毕业设计(论文)是实现人才培养目标的重要实践性环节，对巩固、深化和升华学生所学理论知识，培养学生创新能力、独立工作能力、分析和解决问题能力、工程实践能力起着重要作用。毕业设计(论文)同时也是记录科研成果的重要文献资料，是申请学位的基本依据。为了规范我校本科生毕业设计(论文)的管理，保证毕业设计(论文)质量，特制定《中国石油大学（北京）本科生毕业设计（论文）写作指南》。

\section{一、本科生毕业设计(论文)应体现的三个规范}

本科生毕业设计(论文)的形式与格式虽然不能直接反映论文的学术水平，但可体现出作者的学术修养和文化修养。合格的本科生毕业设计(论文)在形式与格式上体现如下三方面规范：

\begin{enumerate}[label=\arabic*.]
  \item 国家学位条例和学校对学位论文管理的各项具体规范。
  \item 学术著述的一般规范。
  \item 所有正式出版物共同遵循的文字印刷规范。
\end{enumerate}
\section{二、本科生毕业设计(论文)和摘要篇幅要求}

\begin{enumerate}[label=\arabic*.]
  \item 本科生毕业设计(论文)一般为1.5万字左右。
  \item 中文、英文摘要一般为300-500字。
\end{enumerate}
\section{三、本科生毕业设计(论文)封面书写要求}

论文封面包括以下9项：

\begin{enumerate}[label=（\arabic*）]
  \item 单位代码    中国石油大学（北京）的代码是11414；
  \item 学号        以《本科生名册》中的本科生学号为准；
  \item 论文标志    “中国石油大学本科生毕业设计(论文)”；
  \item 论文题目    应简短、醒目、切题，不宜超过25个字；
  \item 学院名称以《本科生培养计划》中的学院为准；
  \item 专业名称以《本科生培养计划》中的专业名称为准；
  \item 学生姓名    以《本科生名册》中的姓名为准；
  \item 指导教师    论文的指导教师；
  \item 论文完成时间    指论文起止时间，用阿拉伯数字书写。
\end{enumerate}
\section{四、本科生毕业设计(论文)摘要书写内容及注意事项}

\begin{enumerate}[label=\arabic*.]
  \item 中文摘要
论文摘要是反映论文内容梗概和研究深度的短文，其目的是使论文评阅人、答辩委员以及读者在较短的时间内对论文有一个简单而又清晰的全面了解。大多数读者看论文时一般首先翻阅论文摘要和目录，因此论文摘要写得好坏，直接决定读者对论文的第一印象，论文作者应对其给予高度重视。

中文摘要字数限制在300-500字左右。为便于检索，要在摘要之后另起一行写出论文的3－5个关键词。前用“关键词”三字加冒号标出，关键词之间用分号隔开，最后一个词后不用标点。

中文摘要在写法上一般不分段落，常采用无人称句。摘要中一般不宜有数学表达式、化学反应式和图表等，不能出现非通用性的外文缩略语或代号，不得引参考文献。

在写作时应注意能反映出以下几方面的内容：

\end{enumerate}
\begin{enumerate}[label=（\arabic*）]
  \item 论文所研究的问题及其目的和意义；
  \item 论文的基本思路和逻辑结构；
  \item 论文的主要方法、内容、结果和结论。
\end{enumerate}
\begin{enumerate}[label=\arabic*.]
  \item 英文摘要
英文摘要应有英文的论文题目。英文摘要的内容必须与中文的内容相对应，但也不是逐字翻译，英文摘要中不能出现汉语的标点符号。

英文摘要应单独从一页开始。

\end{enumerate}
\section{一、本科生毕业设计(论文)应体现的三个规范}

本科生毕业设计(论文)的形式与格式虽然不能直接反映论文的学术水平，但可体现出作者的学术修养和文化修养。合格的本科生毕业设计(论文)在形式与格式上体现如下三方面规范：

\begin{enumerate}[label=\arabic*.]
  \item 国家学位条例和学校对学位论文管理的各项具体规范。
  \item 学术著述的一般规范。
  \item 所有正式出版物共同遵循的文字印刷规范。
\end{enumerate}
\section{二、本科生毕业设计(论文)和摘要篇幅要求}

\begin{enumerate}[label=\arabic*.]
  \item 本科生毕业设计(论文)一般为1.5万字左右。
  \item 中文、英文摘要一般为300-500字。
\end{enumerate}
\section{三、本科生毕业设计(论文)封面书写要求}

论文封面包括以下9项：

\begin{enumerate}[label=（\arabic*）]
  \item 单位代码    中国石油大学（北京）的代码是11414；
  \item 学号        以《本科生名册》中的本科生学号为准；
  \item 论文标志    “中国石油大学本科生毕业设计(论文)”；
  \item 论文题目    应简短、醒目、切题，不宜超过25个字；
  \item 学院名称以《本科生培养计划》中的学院为准；
  \item 专业名称以《本科生培养计划》中的专业名称为准；
  \item 学生姓名    以《本科生名册》中的姓名为准；
  \item 指导教师    论文的指导教师；
  \item 论文完成时间    指论文起止时间，用阿拉伯数字书写。
\end{enumerate}
\section{四、本科生毕业设计(论文)摘要书写内容及注意事项}

\begin{enumerate}[label=\arabic*.]
  \item 中文摘要
论文摘要是反映论文内容梗概和研究深度的短文，其目的是使论文评阅人、答辩委员以及读者在较短的时间内对论文有一个简单而又清晰的全面了解。大多数读者看论文时一般首先翻阅论文摘要和目录，因此论文摘要写得好坏，直接决定读者对论文的第一印象，论文作者应对其给予高度重视。

中文摘要字数限制在300-500字左右。为便于检索，要在摘要之后另起一行写出论文的3－5个关键词。前用“关键词”三字加冒号标出，关键词之间用分号隔开，最后一个词后不用标点。

中文摘要在写法上一般不分段落，常采用无人称句。摘要中一般不宜有数学表达式、化学反应式和图表等，不能出现非通用性的外文缩略语或代号，不得引参考文献。

在写作时应注意能反映出以下几方面的内容：

\end{enumerate}
\begin{enumerate}[label=（\arabic*）]
  \item 论文所研究的问题及其目的和意义；
  \item 论文的基本思路和逻辑结构；
  \item 论文的主要方法、内容、结果和结论。
\end{enumerate}
\begin{enumerate}[label=\arabic*.]
  \item 英文摘要
英文摘要应有英文的论文题目。英文摘要的内容必须与中文的内容相对应，但也不是逐字翻译，英文摘要中不能出现汉语的标点符号。

英文摘要应单独从一页开始。注意事项

  \item 中文摘要
论文摘要是反映论文内容梗概和研究深度的短文，其目的是使论文评阅人、答辩委员以及读者在较短的时间内对论文有一个简单而又清晰的全面了解。大多数读者看论文时一般首先翻阅论文摘要和目录，因此论文摘要写得好坏，直接决定读者对论文的第一印象，论文作者应对其给予高度重视。

中文摘要字数限制在300-500字左右。为便于检索，要在摘要之后另起一行写出论文的3－5个关键词。前用“关键词”三字加冒号标出，关键词之间用分号隔开，最后一个词后不用标点。

中文摘要在写法上一般不分段落，常采用无人称句。摘要中一般不宜有数学表达式、化学反应式和图表等，不能出现非通用性的外文缩略语或代号，不得引参考文献。

在写作时应注意能反映出以下几方面的内容：

\end{enumerate}
\begin{enumerate}[label=（\arabic*）]
  \item 论文所研究的问题及其目的和意义；
  \item 论文的基本思路和逻辑结构；
  \item 论文的主要方法、内容、结果和结论。
\end{enumerate}
\begin{enumerate}[label=\arabic*.]
  \item 英文摘要
英文摘要应有英文的论文题目。英文摘要的内容必须与中文的内容相对应，但也不是逐字翻译，英文摘要中不能出现汉语的标点符号。

英文摘要应单独从一页开始。设计(论文)”；

\end{enumerate}
\begin{enumerate}[label=（\arabic*）]
  \item 论文题目    应简短、醒目、切题，不宜超过25个字；
  \item 学院名称以《本科生培养计划》中的学院为准；
  \item 专业名称以《本科生培养计划》中的专业名称为准；
  \item 学生姓名    以《本科生名册》中的姓名为准；
  \item 指导教师    论文的指导教师；
  \item 论文完成时间    指论文起止时间，用阿拉伯数字书写。
\end{enumerate}
\section{四、本科生毕业设计(论文)摘要书写内容及注意事项}

\begin{enumerate}[label=\arabic*.]
  \item 中文摘要
论文摘要是反映论文内容梗概和研究深度的短文，其目的是使论文评阅人、答辩委员以及读者在较短的时间内对论文有一个简单而又清晰的全面了解。大多数读者看论文时一般首先翻阅论文摘要和目录，因此论文摘要写得好坏，直接决定读者对论文的第一印象，论文作者应对其给予高度重视。

中文摘要字数限制在300-500字左右。为便于检索，要在摘要之后另起一行写出论文的3－5个关键词。前用“关键词”三字加冒号标出，关键词之间用分号隔开，最后一个词后不用标点。

中文摘要在写法上一般不分段落，常采用无人称句。摘要中一般不宜有数学表达式、化学反应式和图表等，不能出现非通用性的外文缩略语或代号，不得引参考文献。

在写作时应注意能反映出以下几方面的内容：

\end{enumerate}
\begin{enumerate}[label=（\arabic*）]
  \item 论文所研究的问题及其目的和意义；
  \item 论文的基本思路和逻辑结构；
  \item 论文的主要方法、内容、结果和结论。
\end{enumerate}
\begin{enumerate}[label=\arabic*.]
  \item 英文摘要
英文摘要应有英文的论文题目。英文摘要的内容必须与中文的内容相对应，但也不是逐字翻译，英文摘要中不能出现汉语的标点符号。

英文摘要应单独从一页开始。

\end{enumerate}
\section{五、本科生毕业设计(论文)目录书写格式及注意事项}

本科生毕业设计(论文)目录是论文的大纲，反映论文的梗概，应将论文目录层次编排清楚。论文的章节按混合编排格式，即：一级标题用“第1章”、“第2章”、……。从二级标题开始，用阿拉伯数字连续编号，在不同层次的数字之间加一个下圆点相隔，最末数字后不加标点。例如：

第一级标题  第1章

第二级标题  1.1

第三级标题  1.1.1

目录中的标题一般编排到三级，正文中的标题层次一般不应超过四级，四级以后可单独编号，如编写作(1) (2) (3) ……或① ② ③……或a. b. c.……等。

本科生毕业设计(论文)总体框架顺序依次为：

\begin{enumerate}[label=（\arabic*）]
  \item 中文封面
  \item 论文独创性声明、论文版权使用授权书（放同一页）
  \item 中文摘要
  \item 英文摘要
  \item 目录
  \item 前言
  \item 正文（第一章至最后一章结论）
  \item 主要符号表
  \item 参考文献
  \item 附录
  \item 致谢
注意事项：

\end{enumerate}
\begin{enumerate}[label=\arabic*.]
  \item 摘要和英文摘要  各单独从一页开始。
  \item 参考文献  只列作者直接阅读过、在正文中被引用过的文献资料。参考文献一律放在论文结论后，不得放在各章之后。每条文献的项目必须完整，诸项缺一不可。各类文献的书写格式均应符合《GB/T 7714-2005文后参考文献著录规则》国家标准。
  \item 致谢  致谢对象限于在学术方面对论文的完成有较重要帮助的人士和团体，感情要诚挚，限300字。
\end{enumerate}
\section{六、本科生毕业设计(论文)前言书写内容及注意事项}

前言在论文正文前。是论文评阅人、答辩委员和读者了解论文研究背景和概貌的主要篇章。主要目的是向论文评阅人、答辩委员和读者阐述论文中所要研究的问题以及与其有关的背景或对一些事项进行说明。

前言通常应包括以下四个方面：

\begin{enumerate}[label=（\arabic*）]
  \item 论文所研究的问题、国内外研究现状以及研究目的和意义；
  \item 论文使用的理论工具、研究方法及技术路线；
  \item 论文的基本思路和逻辑结构；
  \item 论文参考的文献资料、使用的符号、计算公式等需要说明的问题。
前言在写法上不分章节，提倡无人称句。

\end{enumerate}
\section{七、本科生毕业设计(论文)正文书写内容及注意事项}

从第一章开始是学位论文的主体，应注意以下几个问题：

\begin{enumerate}[label=\arabic*.]
  \item 层次要清楚、标题要突出重点,各章之间联系密切，形成一个整体。
  \item 已经被普遍接受而成为常识和学术界公认的观点不要从头叙述。
  \item 如出现一个非通用性的新名词、新术语或新概念，需解释清楚。
  \item 论文中的图、表、附注、公式一律采用阿拉伯数字分章或连续编号。如果图、表中有附注，应采用英文小写字母顺序编号，附注写在图、表的下方。
论文中的插图应具有鲜明性，切忌与列表及文字表述重复。插图中的术语、符号、单位等应同正文表述所用保持一致。插图的序号及图名居中置于图的下方。插图要清楚，坐标比例不要过分放大，同一幅图上不同曲线的点要分别用不同形状标出。

列表中的参数应标明量和单位的符号。列表的序号及列表的名称应置于列表的上方。

公式的编号应加括号，并写在公式的末端，其间不加虚线。

  \item 论文中的量和单位要严格执行GB3100～3102：93有关量和单位的规定。具体要求参阅《常用量和单位》，计量出版社，1996。论文中的单位名称的书写可以采用国际通用符号，也可以用中文名称，但全文应统一，不要两种混用。
  \item 论文中所引用的文献用阿拉伯数字连续标明编号。编号加方括号。
  \item 论文的最后一章是论文的结论，文字必须简明扼要,不要简单重复罗列实验结果。论点应明确、严谨、完整。在写法上不分章节，提倡无人称句。
\end{enumerate}
\section{八、本科生毕业设计(论文)参考文献、附录书写内容及注意事项}

文后参考文献书写采用顺序编码制，即按学位论文中引用参考文献的先后顺序连续编码。本着简单明了，易于查找的原则，根据GB/T 7714-2005 “文后参考文献书写规则”，参照下述格式书写。

\begin{enumerate}[label=\arabic*.]
  \item 专著(注意应标明出版地及所参阅内容在原文献中的位置)，表示方法为：
[序号] 作者.专著名.版本.出版地:出版者,出版年:起止页.

  \item 期刊中析出的文献(注明应标明年、卷（期），尤其注意区分卷和期)，表示方法为：
[序号] 作者.题(篇)名.刊名,出版年,卷号(期号):起止页.

  \item 会议论文，表示方法为：
[序号] 作者.篇名.会议名,会址,开会年:起止页.

  \item 专著(文集)中析出的文献，表示方法为：
[序号] 作者.篇名//文集的编(著)者.文集名.版本.出版地:出版者,出版年:起止页.

  \item 学位论文，表示方法为：
[序号] 作者.题(篇)名:(博(硕)士学位论文).授学位地:授学位单位,授学位年.

  \item 专利文献，表示方法为：
[序号] 专利申请者.专利题名:专利国别,专利号.公告日期.

注意事项：

  \item 每条文献序号用阿拉伯数字按顺序排列，并应与论文中标注的序号相一致；文献序号用[]括住，后加一个空格,不加圆点“.”。
  \item 作者姓名最多书写3个，作者之间以“，”分隔，其后加“等”或“et al.”，最后一名作者后接“.”。
  \item 版本用阿拉伯数字或其他标识，如第2版或2nd ed., 第3版或3rded.,第一版不用书写。
例如：第3版  （原题：第三版）

5th ed.  （原题：Fifth edition）

Rev. ed. （原题：Revised edition）

1998 ed.  (原题：1998 edition)

  \item 出版地只书写出版者所在的城市名称，后接“：”和出版社。
  \item 每条文献最后可加一个黑圆点“．”或者什么也不用加。
  \item 一条文献的各个部分用“.”、“，”、“：”分隔，注意分隔符号的使用位置。
  \item 不要在一条文献段落的中间换页。
附录一般与论文全文装订在一起，如果附录内容很多时可独立成册。

\end{enumerate}
\section{九、本科生毕业设计(论文)书写格式及印刷要求}

\begin{enumerate}[label=\arabic*.]
  \item 论文书写格式
\end{enumerate}
\begin{enumerate}[label=（\arabic*）]
  \item 从摘要开始有页眉，奇数页页眉：黑体，五号，居中。填写内容为“中国石油大学（北京）本科生毕业设计(论文)”。偶数页页眉：黑体，五号，居中。填写内容为“章节标题”。
页脚，标注页码，Arial，五号，居中。

  \item 页码从第一章开始按阿拉伯数字连续编排，Arial，五号字体，页码位于页面底端居中；前置部分用罗马数字（I、II、…）单独编排，宋体,五号。
  \item 正文一级标题居中书写，其它标题靠正文左边书写。
  \item 正文选用宋体小四号书写，用多倍行距:1.25，每行字数：33－35个。
\end{enumerate}
\begin{enumerate}[label=\arabic*.]
  \item 论文印刷
\end{enumerate}
\begin{enumerate}[label=（\arabic*）]
  \item 页边距：上、下、左、右均为3.0cm。
  \item 页眉、页脚：2.0cm。
本科生毕业设计(论文)一律用A4复印纸，纵向单面印刷，封面选用浅黄色纸张。印刷字迹要清晰、端正，版面要美观大方。

\end{enumerate}
\begin{enumerate}[label=\arabic*.]
  \item 论文装订要求
采用左侧1.0cm装订。

其它要求参见格式样本。


\end{enumerate}
% ===========================================================================
%  后置部分
% ===========================================================================

%% --- 结论 ---
\begin{conclusion}
（请在此处填写论文结论）
\end{conclusion}

%% --- 主要符号表 (如需要) ---
% \begin{symbollist}
%   $\alpha$ & 角度 \\
%   $\beta$  & 系数 \\
% \end{symbollist}

%% --- 参考文献 ---
% \makereference{references}

%% --- 附录 ---
% \begin{theappendix}
%   \chapter{附录标题}
%   附录内容...
% \end{theappendix}

%% --- 致谢 ---
\begin{acknowledgement}
（请在此处填写致谢内容，限300字）
\end{acknowledgement}

\end{document}
